\input{../000_header_year.txt}

\begin{document}

%++++++++++++++++++++++++++++++++++++++++++++++++++++++++++++++++++++++++++++++++
%++++++++++++++++++++++++++++++++++++++++++++++++++++++++++++++++++++++++++++++++
%++++++++++++++++++++++++++++++++++++++++++++++++++++++++++++++++++++++++++++++++
\exheader
%++++++++++++++++++++++++++++++++++++++++++++++++++++++++++++++++++++++++++++++++
%++++++++++++++++++++++++++++++++++++++++++++++++++++++++++++++++++++++++++++++++
%++++++++++++++++++++++++++++++++++++++++++++++++++++++++++++++++++++++++++++++++

%%%%%%%%%%%%%%%%%%%%%%%%%%%%%%%%%%%%%%%%%%%%%
%　3mass系の可観測性を調べる問題
%     m0,m1,m2,k1,k2を指定して可観測性を調べる
%     パラメータは２種類でベクトルで定める
% 
%%%%%%%%%%%%%%%%%%%%%%%%%%%%%%%%%%%%%%%%%%%%%


\input{040_def_file.txt}

%================================


%=================================================
	\def\statev
	{
	\left(\begin{array}{c}
	         y_0 \\
		     y_1 \\
		     y_2 \\
		     \frac{dy_0}{dt} \\
		     \frac{dy_1}{dt} \\
		     \frac{dy_2}{dt}
	\end{array}\right)
	}
%
%$\statev$
%
%=================================================


\noindent
{\bf 可観測性の例題}　演習の３質量系で，出力が$y=y_0$の場合，状態方程式と出力方程式は
\begin{align}
 \frac{d}{dt} \statev
 &=
 \left(\begin{array}{cccccc}
        0 & 0 & 0 & 1 & 0 & 0      \\
        0 & 0 & 0 & 0 & 1 & 0      \\
        0 & 0 & 0 & 0 & 0 & 1      \\
	-\frac{k_1}{m_0} -\frac{k_2}{m_0}&  \frac{k_1}{m_0} &  \frac{k_2}{m_0} & 0 & 0 & 0 \\
	 \frac{k_1}{m_1}                & -\frac{k_1}{m_1} &  0               & 0 & 0 & 0 \\
	 \frac{k_2}{m_2}                &  0               & -\frac{k_2}{m_2} & 0 & 0 & 0 
       \end{array}\right)
 \statev
 +
 \left(\begin{array}{c}
        0 \\
        0 \\
        0 \\
	    \frac{1}{m_0} \\
	    0 \\
	    0 
       \end{array}\right)
        F
\nonumber
\\
y &=
 \left(\begin{array}{cccccc}
         1 & 0 & 0 & 0 & 0 & 0
       \end{array}\right)
 \statev
\nonumber
\end{align}
\vspace*{-1cm}

となる．

次のパラメータについて可観測性行列を求めよ.（可観測の場合と不可観測の場合がある）．

・$m_0=\chrmaa$, $m_1=\chrmbb$, $m_2=\chrmca$, $k_1=\chrkaa$, $k_2=\chrkba$

・$m_0=\chrmab$, $m_1=\chrmbb$, $m_2=\chrmcb$, $k_1=\chrkab$, $k_2=\chrkbb$

（計算用に裏面利用可）

\vfill

\begin{tabular}{ll}
	\begin{tabular}{|l|}
		\hline 
		$m_0=\chrmaa$, $m_1=\chrmbb$, $m_2=\chrmca$, $k_1=\chrkaa$, $k_2=\chrkba$の\\
		可観測性行列$N^T$\\
		\hspace*{8cm} \\[3.5cm]
		可観測性：\\
		\hline
	\end{tabular}
	\begin{tabular}{|l|}
		\hline 
		$m_0=\chrmab$, $m_1=\chrmbb$, $m_2=\chrmcb$, $k_1=\chrkab$, $k_2=\chrkbb$の\\
		可観測性行列$N^T$\\
		\hspace*{8cm} \\[3.5cm]
		可観測性：\\
		\hline
	\end{tabular}
\end{tabular}
\newpage
%--------------------------------------------------------------------------------
\noindent
{\bf 可観測性の例題}

\begin{itemize}
\item $m_0=\chrmaa$, $m_1=\chrmbb$, $m_2=\chrmca$, $k_1=\chrkaa$, $k_2=\chrkba$のとき
\end{itemize}

\begin{align}
 \frac{d}{dt} \statev  
 &=
 \chrAa \statev + \chrBa F
\nonumber
\\
 y 
 &= \chrCa \statev 
\nonumber
\end{align}

可観測性のチェック

\begin{align}
 \left(\begin{array}{l}
        C \\
        CA \\
        CA^2  \\
	CA^3 \\
	CA^4 \\
	CA^5
       \end{array}\right)
 &=
 \left(\begin{array}{l}
        C \\
        CA \\
        (CA)A  \\
	(CA^2)A \\
	(CA^3)A \\
	(CA^4)A
       \end{array}\right)
 =
 \chrOba
\nonumber
\end{align}
不可観測




\newpage
%--------------------------------------------------------------------------------
\noindent
{\bf 可観測性の例題}

\begin{itemize}
\item $m_0=\chrmab$, $m_1=\chrmbb$, $m_2=\chrmcb$, $k_1=\chrkab$, $k_2=\chrkbb$のとき
\end{itemize}

\begin{align}
 \frac{d}{dt} \statev  
 &=
 \chrAb \statev + \chrBb F
\nonumber
\\
 y 
 &= \chrCb \statev 
\nonumber
\end{align}

可観測性のチェック

\begin{align}
 \left(\begin{array}{l}
        C \\
        CA \\
        CA^2  \\
	CA^3 \\
	CA^4 \\
	CA^5
       \end{array}\right)
 &=
 \left(\begin{array}{l}
        C \\
        CA \\
        (CA)A  \\
	(CA^2)A \\
	(CA^3)A \\
	(CA^4)A
       \end{array}\right)
 =
 \chrObb
\nonumber
\end{align}
可観測




\end{document}

